\section{Application Validation: Implementation and Proofs}
\label{app:applicationcertificate}

The following additions deploy the minimum-benefit rule on the Bunny and
housing. This is a separate experiment with new seeds, geometries selected
under an acceptance rule, and retained validation samples. It does not
retroactively certify the earlier ranking experiments. Three elementary
constructions make the validator usable: downward clipping of nonnegative
walk contributions, conservative boundary-distance queries, and explicit
source enclosures. The probability guarantee retains the ideal conditional
sampling assumptions of \cref{prop:onesidedwos}. Numerical containment and
source safeguards alone do not prove those sampling laws for finite-precision
positions and a pseudorandom generator.

\subsection{A sharper certificate from the average cavity trace}

The temperature-floor certificate in \cref{thm:coolingbenefit} is sufficient
but can be conservative near existing cavities or a nonuniform heat source.
The same energy identity admits a stronger use of a center query, without
requiring a pointwise temperature floor on the new cavity.

\begin{proof}[Proof of \cref{thm:averagecertificate}]
\textbf{Step 1: difference energy.}
Let $u$ be the current state and $v$ the state after removing
$K=\overline B_r(z)$, extended by zero. Set $w=u-v$ and extend it by zero
outside $D$. As in \cref{lem:floorcapacity},
\[
 B=J(D)-J(D\setminus K)=\int_{\R^3}|\nabla w|^2.
\]
The trace of $w$ on $\partial K$ is the trace of $u$, since $v$ has zero
trace there.

\textbf{Step 2: pair with the equilibrium potential.}
Define $P(x)=1$ for $|x-z|\le r$ and $P(x)=r/|x-z|$ otherwise. On the
exterior of the ball it is harmonic. Its outward normal derivative for
that exterior domain, on its inner sphere, is $+1/r$. Integration by parts
therefore gives
\begin{equation}
 \int_{\R^3}\nabla w\cdot\nabla P
 =\frac1r\int_{\partial B_r(z)}u\,dS
 =4\pi r\,\overline u_r,
 \label{eq:tracecapacitypairing}
\end{equation}
where $\overline u_r=(4\pi r^2)^{-1}\int_{\partial B_r(z)}u\,dS$.
To justify this calculation for $H^1$ functions, first use smooth compactly
supported approximations to the zero extension of $w$. Integrate on the
exterior annulus with an outer sphere beyond their support, so the outer
boundary term vanishes. Pass to the limit by gradient-norm convergence
and the continuous trace map on an annular neighborhood of the inner
sphere. The potential has zero interior gradient and energy $4\pi r$,
as computed in \cref{lem:floorcapacity}.

\textbf{Step 3: apply Cauchy--Schwarz to the energy pairing.}
Squaring \cref{eq:tracecapacitypairing} and applying Cauchy--Schwarz gives
\[
 (4\pi r\,\overline u_r)^2
 \le\left(\int|\nabla w|^2\right)
       \left(\int|\nabla P|^2\right)=B\,4\pi r.
\]
Division by $4\pi r>0$ proves $B\ge4\pi r\,\overline u_r^2$.

\textbf{Step 4: recover the average trace from a center value.}
The centered ball Green identity gives
\begin{equation}
 u(z)=\overline u_r+I_r,\qquad
 I_r=\int_{B_r(z)}G_{B_r(z)}(z,y)f(y)\,dy.
 \label{eq:tracefromcenter}
\end{equation}
If $I_r\le I_r^+$ and $L_u\le u(z)$, then
$\overline u_r\ge L_u-I_r^+$. Also $\overline u_r\ge0$ because $u\ge0$.
Thus $\overline u_r\ge(L_u-I_r^+)_+$, and substitution into Step 3 proves
the claimed bound. Positivity of $G$ and $f\le f_+$ on the inserted ball
give $I_r^+=f_+r^2/6$, using $\int G=r^2/6$.

\textbf{Step 5: affine sources and uniformity.}
For $f(y)=f(z)+a\cdot(y-z)$, the integral of $G(z,y)a\cdot(y-z)$ vanishes:
reflection $y-z\mapsto-(y-z)$ preserves the ball and the radial kernel and
reverses this integrand's sign. Hence $I_r=f(z)r^2/6$ exactly.
All preceding inequalities hold simultaneously for every contained radius
on the one event $\{L_u\le u(z)\}$. They also hold together with the
temperature-floor certificate, so their maximum is a valid benefit lower
bound on that event. The acceptance and adaptive-selection conclusions
follow from the same event-containment proof as \cref{prop:adaptivebenefit};
no additional confidence allocation is needed to compare deterministic
certificates built from the same covered state bound.
\end{proof}

This argument uses classical equilibrium-potential duality and the centered
Green representation. Its role here is a sharper finite-insertion validator,
not a claim that the underlying capacity identity is new. It avoids the
harmonic minimum factor in \cref{eq:temperaturefloor}; the local source
upper bound now concerns the inserted ball itself. The earlier floor and
source-only bounds remain useful when the state estimate is too weak.

\subsection{Clipping gives a smaller, known sample range}

\begin{proposition}[Downward-clipped validation walks]
\label{prop:clippedvalidation}
Under the one-sided construction of \cref{prop:onesidedwos}, fix $C>0$
before drawing a validation batch. Source samples may use a measurable
lower source $0\le\widetilde f\le f$. Stop by a fixed step count $T$,
possibly earlier when the accumulated nonnegative contribution reaches $C$,
and output $W_C=\min(C,\sum_{k<\tau}S_k)$. Then
\begin{equation}
 0\le W_C\le C,\qquad \E W_C\le u(z).
 \label{eq:clippedvalidation}
\end{equation}
For $N$ independent complete walks, the confidence rule in
\cref{eq:onesidedwos} holds with $A=C$.
\end{proposition}
\begin{proof}
\textbf{Step 1: lower sources preserve a one-step inequality.}
Positivity of the ball Green kernel gives
$\int G\widetilde f\le\int Gf$. The one-ball representation therefore
implies
\[
 \E[S_k+u(X_{k+1})\mid\mathcal G_k]\le u(X_k)
\]
at every active step. After stopping, pad with $S_k=0$ and
$X_{k+1}=X_k$; the inequality remains true. Reaching the cap after a step
is a permissible stopping rule because the decision uses only the history
available at the next step.

\textbf{Step 2: telescope and then clip.}
Summing the $T$ finite conditional inequalities and taking expectations yields
$\E\sum_{k<\tau}S_k+\E u(X_\tau)\le u(z)$.
Since $u\ge0$ and
$0\le W_C\le\sum_{k<\tau}S_k$, it follows that $\E W_C\le u(z)$.
Clipping also ensures $W_C\le C$ pathwise. The bound makes no assertion
that clipping preserves the mean; additional downward bias is expected.

\textbf{Step 3: apply bounded-variable concentration.}
The moment-generating-function proof in \cref{prop:onesidedwos}, Step 3,
uses only independence and the range width. Substituting $C$ for $A$ gives
\[
 \Prob\!\left\{\overline W_C-\E W_C>
 C\sqrt{\frac{\log(1/\delta)}{2N}}\right\}\le\delta.
\]
Subtract this allowance and clip the resulting state lower bound at zero.
Step 2 and $u(z)\ge0$ prove its coverage. A smaller $C$ improves the
sampling allowance but can increase clipping bias, so it need not improve
the final benefit certificate.
\end{proof}

\subsection{A conservative distance query on triangle geometry}

\begin{proposition}[Bounding-box distance and contained steps]
\label{prop:boxdistance}
Let a closed embedded triangle surface $S=\bigcup_i T_i$ bound the material
domain, and let $Q_i\supset T_i$ be the coordinate-aligned bounding box of
triangle $i$. For any point $x$,
\begin{equation}
 d_Q(x):=\min_i\operatorname{dist}(x,Q_i)
 \le\operatorname{dist}(x,S).
 \label{eq:boxdistance}
\end{equation}
If $x$ is inside the current material domain, subtracting retained spherical
cavities gives the valid distance lower bound
\[
 d_*(x)=\min\left(d_Q(x),\min_j(|x-c_j|-r_j)\right).
\]
For $0<s<d_*(x)$, the closed ball $\overline B_s(x)$ lies strictly inside
the material domain. A computed endpoint $x'$ is still interior whenever
a certified upper bound on $|x'-x|$ is smaller than a certified lower bound
on $d_*(x)$.
\end{proposition}
\begin{proof}
\textbf{Step 1: set containment reverses distance.}
For each $i$, the infimum of $|x-y|$ over the larger set $Q_i$ is at most
its infimum over $T_i$. Taking minima over $i$ proves
\cref{eq:boxdistance}. A box may reach into the material and give zero
distance at an interior point; stopping there is conservative.

\textbf{Step 2: evaluate the box distance.}
If $Q=\prod_{a=1}^3[l_a,h_a]$, its nearest point is the coordinatewise
clamping of $x_a$ to $[l_a,h_a]$. Indeed, the squared Euclidean distance is
a sum of three independent nonnegative coordinate objectives. Thus
\[
 \operatorname{dist}(x,Q)^2
 =\sum_{a=1}^3\max(l_a-x_a,x_a-h_a,0)^2.
\]
Lower bounds on each nonnegative coordinate gap, their squares, their sum,
and its square root give a lower bound on the distance. The implementation
rounds each of these elementary operations outward using adjacent binary64
numbers, without fast-math reassociation. This statement assumes the usual
correctly rounded IEEE arithmetic and square root for those operations.

\textbf{Step 3: justify hierarchy pruning.}
Store in every internal hierarchy node a box containing all descendant
triangle boxes. Its lower distance is a lower bound on the distance to each
descendant triangle. If the current returned candidate lower bound is $b$
and a node's lower distance is at least $b$, every triangle in that node
has true distance at least $b$, so pruning that node cannot make $b$ exceed
the global true surface distance. Evaluating all unpruned leaves and taking
their minimum lower bound preserves this invariant. Using lower bounds
rather than nearest-point estimates is essential to this argument.

\textbf{Step 4: include old cavities and strict containment.}
For $x$ outside $\overline B_{r_j}(c_j)$, its distance to that sphere is
$|x-c_j|-r_j$. Taking the minimum with the outer-surface lower bound gives
a lower bound for the distance to every boundary of the current domain.
If a point of $\overline B_s(x)$ left the domain, the line segment from
interior $x$ to that point would first cross a boundary at distance at most
$s<d_*(x)$, a contradiction. The same segment argument with
$|x'-x|<d_*(x)$ proves the computed-endpoint claim.
\end{proof}

The initial finalist is classified by exact rational positive-axis ray
parity, treating stored binary64 mesh coordinates as exact dyadic rationals.
Projected bounding boxes cull triangles; the remaining projected orientation
and intersection-sign tests use rational arithmetic. On an edge or coplanar
hit the code retries another axis and rejects the classification if all axes
are degenerate. For a closed embedded surface and a nondegenerate ray, each
crossing changes interior status; the ray ends outside the bounded object,
so an odd crossing count is equivalent to an interior origin. Prior cavity
exclusion is checked by exact squared-distance comparison. These predicates
presuppose the closed embedded input surface; manifold diagnostics are not
an exhaustive proof of absence of every self-intersection.

After initialization, every source location and walk endpoint passes the
upper-distance guard in \cref{prop:boxdistance}. A guard failure invalidates
the validation record; it is not silently counted as a zero-valued sample.
Steps use $s=\min(s_*,0.95\,\underline d_*)$, and the source integration
and positive accumulation are rounded downward. A rounded endpoint remaining
inside is a geometric guarantee; it does not prove that its distribution
is exactly the ideal uniform spherical measure.

\subsection{Explicit local source bounds}

\begin{proposition}[Affine and Gaussian source enclosures]
\label{prop:sourcesenclosed}
For $f(x)=1+a(x_2+1/2)$ with $a\ge0$, valid extrema on $B_R(z)$ are
$f_\pm=1+a(z_2+1/2\pm R)$. For a nonnegative Gaussian mixture
\[
 f(x)=b+\sum_j a_j\exp\!\left(-\sum_{k=1}^3
                     \frac{(x_k-c_{jk})^2}{w_{jk}^2}\right),
\]
assume $b,a_j\ge0$ and $w_{jk}>0$.
define
\begin{align*}
 t_j^{\min}&=\sum_k\frac{\max(|z_k-c_{jk}|-R,0)^2}{w_{jk}^2},\\
 t_j^{\max}&=\sum_k\frac{(|z_k-c_{jk}|+R)^2}{w_{jk}^2}.
\end{align*}
Then $b+\sum_j a_j e^{-t_j^{\max}}$ and
$b+\sum_j a_j e^{-t_j^{\min}}$ are respectively lower and upper source
bounds throughout $B_R(z)$. Moreover, for $t\ge0$ and integer $n\ge1$,
\begin{equation}
 \max(0,1-t/n)^n\le e^{-t}\le(1+t/n)^{-n}.
 \label{eq:rationalexpenclosure}
\end{equation}
\end{proposition}
\begin{proof}
\textbf{Step 1: enclose coordinates.}
Every point in $B_R(z)$ has $z_k-R\le x_k\le z_k+R$.
Monotonicity of the affine source gives its two bounds. The triangle and
reverse-triangle inequalities give
\[
 \max(|z_k-c_{jk}|-R,0)\le|x_k-c_{jk}|
                  \le|z_k-c_{jk}|+R.
\]
Squaring, dividing by $w_{jk}^2>0$, and summing proves that the exponent
argument lies in $[t_j^{\min},t_j^{\max}]$.

\textbf{Step 2: enclose each exponential.}
The function $t\mapsto e^{-t}$ is decreasing. Apply it to the exponent
interval, multiply by $a_j\ge0$, and sum with $b$ to obtain the source
enclosures. Optimizing each coordinate separately relaxes the ball to a box;
this weakens the bounds but does not invalidate them.

\textbf{Step 3: remove dependence on a rounded exponential function.}
For $0\le a<1$, integration of $1/(1-s)\ge1$ over $[0,a]$ gives
$\log(1-a)\le-a$. With $a=t/n<1$, multiply by $n$ and exponentiate to
get the lower inequality in \cref{eq:rationalexpenclosure}; for $t\ge n$
its left side is zero and the result is immediate. Similarly, integration
of $1/(1+s)\le1$ gives $\log(1+a)\le a$ for $a\ge0$. Negate, multiply
by $n$, and exponentiate to get $(1+t/n)^{-n}\ge e^{-t}$.
The implementation uses $n=1024$, computed by ten squarings. Downward
rounding in the lower formula and a lower denominator followed by upward
reciprocation in the upper formula preserve the enclosures. For a sampled
source point, only the nonnegative lower formula is used. It supplies the
lower-source condition in \cref{prop:clippedvalidation}.
\end{proof}

\subsection{Predeclared protocol and scope of its probability statement}

There are six attempted insertions per object. Each attempt scouts the
current retained material, rescores a shortlist, and fixes one finalist
before collecting any certificate samples. The shortlist's scores are
used only for selection; they supply no certificate. A separate validation
stream evaluates that finalist at the four nested sample sizes
$N\in\{1024,4096,16384,65536\}$. We allocate
$\delta_{tj\ell}=1/4800$ to each of the $2\cdot6\cdot4=48$ possible
tests. Their total is exactly $0.01$. \Cref{prop:adaptivebenefit} therefore
covers the maximum of the four bounds, selection of the containing radius,
and all accepted updates, conditional on the stated oracle assumptions.

For a conservative numerical concentration allowance, use
\[
 a_N=C\sqrt{8.5/(2N)}
 \ge C\sqrt{\log(4800)/(2N)}.
\]
The constant needs no numerical logarithm: the positive exponential series
gives
$e^{17/2}>\sum_{k=0}^{25}(17/2)^k/k!>4800$, where the last comparison is
checked as an exact rational inequality.
Monotonicity of $\log$ proves $\log4800<8.5$.
Sample means are recomputed by exact dyadic summation of the stored
binary64 values, then rounded down. The allowance is rounded up, and the
subsequent benefit expression is evaluated with lower arithmetic bounds.
The $\pi$ lower constant is obtained from Machin's identity
$\pi=16\arctan(1/5)-4\arctan(1/239)$ and alternating rational partial sums
(12 terms for the first arctangent, 5 for the second). The alternating-series
error signs give a lower first sum and upper second sum, hence a lower
$\pi$. For completeness, the tangent double-angle formula gives
$\tan(2\arctan(1/5))=5/12$ and
$\tan(4\arctan(1/5))=120/119$. The tangent subtraction formula then gives
\[\tan(4\arctan(1/5)-\arctan(1/239))=1.\] This angle lies in $(0,\pi/2)$,
so it equals $\pi/4$, proving the stated identity. These scalar arithmetic precautions do not bound the error in the
random walk's ideal sampling distribution.

The primary targets are fixed absolute compliance reductions:
$2\times10^{-6}$ for the Bunny and $1.5\times10^{-5}$ for the housing.
They equal $0.1\%$ of stipulated reference scales $0.002$ and $0.015$,
respectively. These scales are fixed design units, not confidence bounds on
the true baseline compliance. Additional $0.5\%,1\%,2\%$ thresholds are
diagnostic tests on the same proposals and current domains, not separately
optimized designs. A percentage of a noisy measured baseline is never used
as a guaranteed relative target.

Only accepted cavities enter the next iteration's material domain. For
each proposed insertion, separate objective evaluations compare its current
and trial domains with four randomized spatial blocks, 512 original-domain
probes per block, and 256 WoS paths per active probe. Their differences
estimate the actual finite benefit and its block standard error; these
approximate references do not participate in acceptance. The original-domain
volume normalization includes zero values inside removed cavities.

All four validation budgets are computed for a tightness diagnostic, even
when an earlier budget already passes. We report that complete diagnostic
cost. The first passing stage and its prefix time describe the smaller
computation that a staged stopping policy could use; they do not erase the
extra work actually executed. Selection, validation, hierarchy construction,
and independent reference evaluation are kept distinct.

\subsection{Measured validation results}
\label{app:applicationvalidationresults}

\Cref{tab:applicationvalidation} reports a separate six-attempt design on
each object. Eleven insertions pass the fixed absolute target at the
first predeclared stage of 1,024 walks; the fifth housing insertion first
passes at 4,096 walks. Later samples tighten the reported
benefit bounds; they are not required to make these particular accept/reject
decisions. The maximum of the four covered stage bounds supplies each
reported lower benefit. Its sum over accepted insertions is justified by
\cref{prop:adaptivebenefit}. Displayed cumulative lower bounds are rounded
down, and the raw samples and full-precision bounds are retained.

\begin{table}[t]
\centering\small
\caption{New guarded designs, with one finalist per attempt and a total
ideal-sampling error allocation of $0.01$ across both objects and all
48 validation stages. Compliance is unnormalized. The two numerical
references are approximations, not exact ground truth. Times cover the
respective kernels over six attempts; reference evaluation, process startup,
and hierarchy construction are excluded.}
\label{tab:applicationvalidation}
\begin{tabular}{lrr}
\toprule
Quantity & Bunny & Housing\\
\midrule
Accepted / attempted & $6/6$ & $6/6$\\
Target per insertion & $2\cdot10^{-6}$ & $1.5\cdot10^{-5}$\\
Cumulative lower benefit & $5.75\cdot10^{-4}$ & $6.21\cdot10^{-3}$\\
WoS reference benefit & $8.448\cdot10^{-4}$ & $1.020\cdot10^{-2}$\\
Multigrid reference benefit & $8.455\cdot10^{-4}$ & $1.013\cdot10^{-2}$\\
Scout + rescore (s) & $2.97$ & $0.53$\\
First passing prefixes (s) & $2.82$ & $1.48$\\
Full validation sweep (s) & $150.39$ & $52.17$\\
\bottomrule
\end{tabular}
\end{table}

The new designs reduce the independently estimated objective by $47.18\%$
on the Bunny and $68.65\%$ on the housing. These measured percentages use
their respective new reference baselines and are not themselves certified
relative reductions. The full-precision cumulative bounds are $68.1\%$
and $61.0\%$ of the corresponding estimated absolute benefits. Per insertion,
the average-trace bound is $2.4$--$16.9\times$ the temperature-floor bound
on the Bunny and $59.4$--$365.4\times$ on the housing, evaluated on the
same validation data and current domains. This compares two derived
certificates, not two forward-solver runtimes or an external state-of-the-art
optimization baseline.

For an additional deterministic check, geometric multigrid uses grids of
$97\times113\times85$ on the Bunny and $113\times73\times57$ on the
housing, with relative residual tolerance $10^{-8}$. Its positive state
values are recovered from the solver's squared-state score output, and
objectives use the same independent spatial probes as the WoS reference.
Its total-benefit estimates differ from WoS by $0.085\%$ and $0.645\%$,
respectively. Every reported insertion bound lies below both approximate
reference estimates. This agreement is a numerical consistency check; it
does not bound either method's continuum error or establish exact coverage.

The complete diagnostic collects $786{,}432$ independent-walk records under
the statistical model and executes $45{,}470{,}847$ guarded steps with no
containment failure. The audit recomputes all 48 state bounds from retained
samples, checks the twelve initial classifications with exact rational
predicates, compares conservative center distances with libigl, and tests
240 source evaluations using high-precision arithmetic. The geometric and
source inequalities are justified by \cref{prop:boxdistance,prop:sourcesenclosed};
the finite audit is an implementation check, not their proof.

The conservative query backend and large sample sweep are expensive. The
first-passing prefix times in the table are recorded kernel times from the
executed stages, illustrating what staged acceptance could save; the complete
sweep actually ran. Neither those prefix times nor the earlier ranking-only
speedups establish that fully validated optimization is faster than the
deterministic baseline. Validation tightness, overhead, and the remaining
ideal-sampling assumption are reported as separate properties.

